% =====================================================================
%  MINI CURSO DO CEM — Introdução à IA Generativa nas Ciências Sociais
%  Formato: 2 dias (8h00–12h30) · introdutório · independente de plataforma
% =====================================================================
\documentclass[11pt,aspectratio=169,t]{beamer}
% --- Codificação e idioma -----------------------------------------------------
\usepackage[utf8]{inputenc}
\usepackage[T1]{fontenc}
\usepackage[brazilian]{babel}
\usepackage{lmodern}
% --- Tema visual --------------------------------------------------------------
\usetheme{Madrid}
\usecolortheme{default}
\useinnertheme{rectangles}
\useoutertheme{miniframes}
\definecolor{azulUNI}{RGB}{31,78,121}
\definecolor{azulUNIclaro}{RGB}{68,114,150}
\definecolor{laranjaDestaque}{RGB}{200,90,30}
\definecolor{verdeCodigo}{RGB}{0,110,80}
\definecolor{cinzaCodigo}{RGB}{245,245,245}
\definecolor{cinzaSaida}{RGB}{237,240,243}
\definecolor{cinzaescuro}{RGB}{58,58,58}
\setbeamercolor{structure}{fg=azulUNI}
\setbeamercolor{frametitle}{bg=azulUNI,fg=white}
\setbeamercolor{title}{bg=azulUNI,fg=white}
\setbeamercolor{section in head/foot}{bg=azulUNIclaro,fg=white}
\setbeamercolor{subsection in head/foot}{bg=azulUNI,fg=white}
\setbeamercolor{author in head/foot}{bg=azulUNI,fg=white}
\setbeamercolor{title in head/foot}{bg=azulUNIclaro,fg=white}
\setbeamercolor{date in head/foot}{bg=azulUNI,fg=white}
\setbeamercolor{block title}{bg=azulUNI,fg=white}
\setbeamercolor{block body}{bg=cinzaSaida,fg=cinzaescuro}
\setbeamercolor{block title alerted}{bg=laranjaDestaque,fg=white}
\setbeamercolor{block body alerted}{bg=cinzaSaida,fg=cinzaescuro}
\setbeamercolor{block title example}{bg=verdeCodigo,fg=white}
\setbeamercolor{block body example}{bg=cinzaSaida,fg=cinzaescuro}
\setbeamercolor{alerted text}{fg=laranjaDestaque}
\setbeamersize{text margin left=0.55cm, text margin right=0.55cm}
% --- Pacotes auxiliares -------------------------------------------------------
\usepackage{tikz}
\usetikzlibrary{shapes.geometric,arrows.meta,positioning,calc,fit,backgrounds}
\usepackage{booktabs}
\usepackage{array}
\usepackage{amsmath,amssymb}
\usepackage{xcolor}
\usepackage{graphicx}
\usepackage{hyperref}
\usepackage{listings}
\usepackage{multicol}
\hypersetup{colorlinks=true,linkcolor=azulUNI,citecolor=laranjaDestaque,urlcolor=azulUNI}
% ---------- Listings (para prompts) ----------
\lstset{
  basicstyle=\ttfamily\scriptsize,
  backgroundcolor=\color{cinzaCodigo},
  frame=single,
  rulecolor=\color{cinzaescuro},
  breaklines=true,
  breakatwhitespace=true,
  columns=fullflexible,
  showstringspaces=false,
  xleftmargin=4pt,
  xrightmargin=4pt,
  aboveskip=4pt,
  belowskip=4pt,
  inputencoding=utf8,
  extendedchars=true,
  literate=
    {á}{{\'a}}1 {é}{{\'e}}1 {í}{{\'i}}1 {ó}{{\'o}}1 {ú}{{\'u}}1
    {Á}{{\'A}}1 {É}{{\'E}}1 {Í}{{\'I}}1 {Ó}{{\'O}}1 {Ú}{{\'U}}1
    {à}{{\`a}}1 {è}{{\`e}}1 {ì}{{\`i}}1 {ò}{{\`o}}1 {ù}{{\`u}}1
    {À}{{\`A}}1 {È}{{\`E}}1 {Ì}{{\`I}}1 {Ò}{{\`O}}1 {Ù}{{\`U}}1
    {â}{{\^a}}1 {ê}{{\^e}}1 {î}{{\^i}}1 {ô}{{\^o}}1 {û}{{\^u}}1
    {Â}{{\^A}}1 {Ê}{{\^E}}1 {Î}{{\^I}}1 {Ô}{{\^O}}1 {Û}{{\^U}}1
    {ã}{{\~a}}1 {õ}{{\~o}}1 {Ã}{{\~A}}1 {Õ}{{\~O}}1
    {ç}{{\c{c}}}1 {Ç}{{\c{C}}}1
    {ü}{{\"u}}1 {Ü}{{\"U}}1
}
% ---------- Slide de transição entre blocos ----------
\newcommand{\transicao}[2]{%
  \begin{frame}[plain]
    \centering
    \vspace*{2.2cm}
    {\large\color{cinzaescuro} #1}\\[0.8cm]
    {\Large\color{azulUNI}\textbf{#2}}
  \end{frame}}
% ---------- Slide de abertura de dia ----------
\newcommand{\aberturadia}[3]{%
  \begin{frame}[plain]
    \centering
    \vspace*{1.6cm}
    {\Huge\color{azulUNI}\textbf{#1}}\\[0.6cm]
    {\large\color{cinzaescuro} #2}\\[1.2cm]
    {\normalsize\color{laranjaDestaque}\textbf{#3}}
  \end{frame}}
% =====================================================================
%  METADADOS
% =====================================================================
\title[Introdução à IAG — Mini curso CEM]{Introdução à Inteligência Artificial Generativa\\
\large nas Ciências Sociais}
\subtitle{Epistemologia, Pesquisa, Ensino e Escrita com IA \\ \small Mini curso do CEM · 2 dias · nível introdutório}
\author[Anderson Henrique]{\textbf{Anderson Henrique}\\
\small CEM/USP $\bullet$ INCT QualiGov\\
\small Pesquisador Convidado IPEA\\
\small Doutor em Ciência Política pela UFPE}
\institute{}
\date{2026}
% =====================================================================
%  DOCUMENTO
% =====================================================================
\begin{document}

\begin{frame}[plain]
  \titlepage
\end{frame}

\begin{frame}{Como este mini curso funciona}
  \begin{block}{Formato}
    Dois dias, 8h00--12h30 $\bullet$ 25 participantes $\bullet$ sala comum, cada um com seu laptop.
  \end{block}
  \pause
  \begin{alertblock}{Nível introdutório}
    Feito para quem \textbf{domina pouco ou nada} da ferramenta. Não há pergunta boba. Ninguém precisa saber programar.
  \end{alertblock}
  \pause
  \begin{exampleblock}{Independente de plataforma}
    Tudo funciona em \textbf{qualquer IA generativa}, como ChatGPT, Gemini ou Claude, entre outras. \textit{Recomendamos o Claude}, mas use a que você já tem.
  \end{exampleblock}
\end{frame}

\begin{frame}{O caminho dos dois dias}
  \footnotesize
  \begin{columns}[t]
    \column{0.5\textwidth}
    \textbf{\color{azulUNI}Dia 1. Entender e decidir}
    \begin{itemize}
      \item M0. Abertura e diagnóstico
      \item M1. O que é IA generativa
      \item M2. Epistemologia da IA generativa
      \item M3. Sete problemas documentados
      \item M4. Ética, integridade e regulação \\ \scriptsize (com debate de casos)
    \end{itemize}
    \column{0.5\textwidth}
    \textbf{\color{azulUNI}Dia 2. Praticar e aplicar}
    \begin{itemize}
      \item M5. Escrevendo bons prompts
      \item M6. IA no ensino: aulas e apresentações
      \item M7. IA na pesquisa: ferramentas e modelos
      \item M8. IA na escrita acadêmica
      \item M9. Desdobramentos contemporâneos: agentes, \textit{skills} e código
      \item M10. Síntese e encerramento
    \end{itemize}
  \end{columns}
  \vfill
  \begin{block}{}
    \centering O Dia 1 responde \textit{``o que é isso e o que posso fazer?''}. O Dia 2 responde \textit{``como faço?''}.
  \end{block}
\end{frame}

% =====================================================================
\section{Dia 1}
\aberturadia{DIA 1}{Entender e decidir}{8h00 -- 12h30}
% =====================================================================
% ---------------- M0 ----------------
\subsection{M0. Abertura e diagnóstico}
\transicao{Módulo 0}{Abertura e diagnóstico}
\begin{frame}{Antes de tudo: quem está na sala?}
  \begin{block}{Diagnóstico rápido...}
    \begin{enumerate}
      \item Quem \textbf{nunca} usou nenhuma IA generativa?
      \item Quem já usou, mas \textbf{sem saber direito} o que estava fazendo?
      \item Quem usa \textbf{com alguma frequência} no trabalho?
    \end{enumerate}
  \end{block}
  \pause
  \vfill
  \begin{exampleblock}{}
    \centering Este curso foi desenhado para os grupos \textbf{1 e 2}. Se você está no 3, vai ganhar \textbf{o critério} (essa parte quase ninguém tem, viu?)
  \end{exampleblock}
\end{frame}
\begin{frame}{Primeiro contato, todos juntos}
  \begin{block}{Abra a IA que você tiver e envie esse prompt}
    \centering
    \textit{``Explique, em um parágrafo, para um calouro de Ciências Sociais, o que é o conceito de capacidade estatal.''}
  \end{block}
  \pause
  \vfill
  \begin{itemize}
    \item Não importa a resposta.
    \item Guarde o resultado. Vamos \textbf{dissecá-lo} daqui a pouco.
    \item \alert{Alguém recebeu algo que parecia bom demais? Ou que citou algo que você não sabe se existe?}
  \end{itemize}
\end{frame}
% ---------------- M1 ----------------
\subsection{M1. O que é IA generativa}
\transicao{Módulo 1}{O que é IA generativa?}
\begin{frame}{Como funciona, previsão e não conhecimento}
  \centering
  \begin{tikzpicture}[node distance=0.5cm and 0.9cm]
    \tikzstyle{caixa}=[rectangle, rounded corners, draw=azulUNI, very thick,
                       fill=cinzaSaida, text width=2.7cm, align=center, minimum height=1.3cm]
    \tikzstyle{saida}=[rectangle, rounded corners, draw=verdeCodigo, very thick,
                       fill=cinzaSaida, text width=2.7cm, align=center, minimum height=1.3cm]
    \tikzstyle{seta}=[-{Latex[length=2.5mm]}, very thick, azulUNI]
    \node[caixa] (entrada) {\textbf{Entrada}\\\scriptsize ``A capital de Pernambuco é''};
    \node[caixa, right=of entrada] (tokens) {\textbf{Tokenização}\\\scriptsize o texto vira sequência numérica};
    \node[caixa, right=of tokens] (modelo) {\textbf{Modelo}\\\scriptsize calcula a probabilidade do próximo token};
    \node[saida, right=of modelo] (saida) {\textbf{Saída}\\\scriptsize ``Recife'' (p $\approx$ 0,97)};
    \draw[seta] (entrada) -- (tokens);
    \draw[seta] (tokens) -- (modelo);
    \draw[seta] (modelo) -- (saida);
  \end{tikzpicture}
  \pause
  \vfill
  \begin{alertblock}{A ideia mais importante do dia}
    O modelo \textbf{não consulta fatos}, mas a continuação mais provável. Quando a continuação mais provável é falsa, ele a gera \textbf{com a mesma confiança}.
  \end{alertblock}
\end{frame}
\begin{frame}{Uma analogia}
  \begin{block}{}
    O corretor automático do seu celular sugere a próxima palavra a partir do que você mais digitou antes. Uma IA generativa faz a mesma aposta, só que em escala industrial, treinada em bilhões de páginas, e a aposta vira um parágrafo inteiro em vez de uma palavra.
  \end{block}
  \pause
  \vfill
  \begin{itemize}
    \item Ela não ``sabe'' que Recife é a capital de Pernambuco;
    \item Ela calculou que, depois daquelas palavras, ``Recife'' é o que quase sempre vem.
    \pause
    \item \alert{Por isso ela erra com confiança. Não existe, do lado de lá, alguém verificando.}
  \end{itemize}
\end{frame}

\begin{frame}{Voltando ao seu primeiro prompt}
  \small
  Pegue a resposta que você guardou no Módulo 0, sobre capacidade estatal.
  \begin{itemize}
    \item Ela tinha a estrutura de algo \textbf{correto}, ou você de fato conferiu cada afirmação?
    \item Havia algum dado, autor ou exemplo que parecia \textbf{bom demais} ou que você não sabe se existe?
  \end{itemize}
  \pause
  \vfill
  \begin{alertblock}{}
    \centering Se pareceu ótimo à primeira vista, foi o mecanismo funcionando exatamente como descrito. Bom texto não é o mesmo que texto correto.
  \end{alertblock}
\end{frame}
\begin{frame}{Vocabulários importantes}
  \small
  \begin{columns}[t]
    \column{0.5\textwidth}
    \begin{description}
      \item[Token] o pedacinho em que o texto é fatiado. Custo e limites se medem nele.
      \item[Prompt] o comando que você escreve. \textbf{É a habilidade central};
      \item[Contexto] tudo que a IA ``enxerga'' naquela conversa. Fora dali, ela não sabe.
      \item[Modelo] a versão específica da IA (GPT, Claude, Gemini\ldots).
    \end{description}
    \column{0.5\textwidth}
    \begin{description}
      \item[Papel/persona] pedir que ela ``aja como'' alguém. Muda o tom, mas \textbf{não} cria conhecimento.
      \item[Temperatura] quão ``criativa'' ou previsível é a resposta.
      \item[Alucinação] informação falsa com cara de verdadeira. \alert{Não é defeito, é o mecanismo.}
    \end{description}
  \end{columns}
\end{frame}
\begin{frame}{As principais ferramentas hoje}
  \scriptsize
  \begin{tabular}{p{2.0cm} p{2.0cm} p{3.6cm} p{3.6cm}}
    \toprule
    \textbf{Ferramenta} & \textbf{Quem faz} & \textbf{Costuma ir bem em} & \textbf{Observação} \\
    \midrule
    \textbf{Claude} & Anthropic (EUA) & Textos longos, escrita cuidadosa, instruções detalhadas & \textit{Recomendação deste curso} \\
    \addlinespace
    \textbf{ChatGPT} & OpenAI (EUA) & Uso geral, ecossistema amplo & O mais difundido \\
    \addlinespace
    \textbf{Gemini} & Google (EUA) & Integração com o ecossistema Google & Bom acesso a documentos \\
    \addlinespace
    \textbf{Sabiá-4} & Maritaca AI (Brasil) & Português brasileiro, inclusive jurídico & Modelo nacional (jan. 2026) \\
    \bottomrule
  \end{tabular}
  \pause
  \vspace{0.2cm}
  \begin{exampleblock}{}
    \centering \textbf{Tudo neste curso funciona em qualquer uma delas.} Use a que você já tem e não troque por minha causa.
  \end{exampleblock}
\end{frame}
\begin{frame}{Por que ela inventa referências?}
  \begin{alertblock}{A mecânica}
    A IA prevê \textbf{como se parece} uma referência sobre o tema (autor plausível, periódico real, título convincente, entre outros), sem buscar nada em base alguma.
  \end{alertblock}
  \pause
  \vfill
  \begin{itemize}
    \item Um estudo registrou referências inventadas em \textbf{69\%} das consultas (Alkaissi \& McFarlane, 2023);
    \item Elas parecem perfeitas: Formato certo, autores que existem e revistas que existem. \textbf{Só o artigo é que não existe;}
    \pause
    \item \alert{Regra de ouro nº 1}: toda referência sugerida por IA se verifica na fonte antes de usar.
  \end{itemize}
\end{frame}
\begin{frame}{Cace uma alucinação (Hands-on 1)}
  \begin{exampleblock}{Na sua IA, peça:}
    \textit{``Liste 10 artigos acadêmicos, com referência completa, sobre [um tema específico da sua pesquisa]. Priorize trabalhos brasileiros dos últimos 5 anos.''}
  \end{exampleblock}
  \vfill
  \begin{enumerate}
    \item Escolha 2 referências da lista;
    \item Procure cada uma no Google Scholar ou no Portal CAPES;
    \item \textbf{Quantas existem de verdade?}
  \end{enumerate}
  \vfill
  \begin{block}{}
    \centering Guarde o resultado, tá? Ele vai voltar no Módulo 4!
  \end{block}
\end{frame}
% ---------------- M2 ----------------
\subsection{M2. Epistemologia}
\transicao{Módulo 2}{Epistemologia da IA generativa}
\begin{frame}{A pergunta que organiza o curso}
  \begin{block}{}
    \centering\large
    A IA Generativa muda \alert{o que} fazemos como pesquisadores\\
    ou apenas \alert{como} fazemos?
  \end{block}
  \pause
  \vfill
  \begin{itemize}
    \item Se muda só o \textit{como}, é ferramenta. Aprende-se e segue-se.
    \pause
    \item Se muda o \textit{quê}, mexe com a natureza do nosso trabalho, e aí precisamos discutir.
  \end{itemize}
\end{frame}
\begin{frame}{Conhecimento explícito e conhecimento tácito}
  \begin{block}{Uma distinção clássica}
    \textbf{Conhecimento explícito} é o que pode ser formalizado, escrito, transferido em texto; é exatamente o que um modelo de linguagem processa bem.\\[4pt]
    \textbf{Conhecimento tácito} é o julgamento construído na prática, no campo, na orientação, na leitura repetida; resiste à formalização exaustiva.
  \end{block}
  \pause
  \vfill
  \begin{alertblock}{}
    A IA generativa lida bem com o explícito. Ela formaliza mal ou nem tenta formalizar, o tácito. Aqui entram o faro metodológico, o julgamento sobre relevância e o senso do que ``não fecha'' numa análise.
  \end{alertblock}
  \pause
  \centering\scriptsize Punziano (2025), \textit{Societies}, propõe a ``epistemologia adaptativa'' para pensar essa reconfiguração nas Ciências Sociais.
\end{frame}
\begin{frame}{Quem tem autoridade epistêmica?}
  \footnotesize
  \begin{block}{Autoridade epistêmica}
    O direito reconhecido de dizer o que conta como conhecimento válido. Tradicionalmente, quem testemunhou, mediu, leu e argumentou.
  \end{block}
  \pause
  \vspace{0.1cm}
  \begin{itemize}
    \item A IA gera texto plausível \textbf{sem esse tipo de \textit{warrant}}. Não testemunhou, não mediu, nem decidiu nada.
    \pause
    \item Boriati (2026), \textit{Societies}: a IA como \textit{ator epistêmico} que coproduz sentido, não só instrumento neutro.
    \pause
    \item Aydeniz (2026), \textit{Sociology Compass}: a autoridade epistêmica se desloca quando quem ``sabe'' deixa de ser claramente humano.
  \end{itemize}
  \pause
  \vspace{0.1cm}
  \begin{alertblock}{}
    \centering Se você aceita a resposta da IA sem julgamento, não delegou uma tarefa. Delegou a autoridade.
  \end{alertblock}
\end{frame}

\begin{frame}{Sua resposta do Módulo 0, outra vez}
  \small
  Volte à resposta sobre capacidade estatal que você guardou.
  \begin{itemize}
    \item Ela ficou no conhecimento \textbf{explícito} (definição, dimensões da literatura), ou tentou entrar em \textbf{tácito} (qual exemplo era \textit{o mais relevante})?
    \item Ao aceitar a resposta, você checou cada afirmação, ou emprestou a ela alguma \textbf{autoridade} que não foi conquistada?
  \end{itemize}
  \pause
  \vfill
  \begin{alertblock}{}
    \centering Não há problema em usar a resposta como ponto de partida. O problema é usá-la como ponto final, sem esse julgamento.
  \end{alertblock}
\end{frame}

\begin{frame}{Fechando o círculo, o quê ou o como?}
  \begin{block}{}
    \centering\large
    A IA muda o \textit{como} quando o julgamento humano permanece no comando.\\
    Muda o \textit{quê} quando esse julgamento é terceirizado sem perceber.
  \end{block}
\end{frame}
% ---------------- M3 ----------------
\subsection{M3. Sete problemas documentados}
\transicao{Módulo 3}{Sete problemas documentados}
\begin{frame}{Sete desafios que precisamos encarar}
  \begin{columns}[t]
    \column{0.5\textwidth}
    \textbf{\color{laranjaDestaque}Problemas da máquina}
    \begin{enumerate}
      \item \textbf{Alucinação}: inventa com confiança;
      \item \textbf{Viés}: reproduz o mundo dos dados em que foi treinada;
      \item \textbf{Opacidade}: nem quem a fez explica a resposta;
      \item \textbf{Irreprodutibilidade}: o mesmo prompt gera respostas diferentes a cada rodada.
    \end{enumerate}
    \column{0.5\textwidth}
    \textbf{\color{laranjaDestaque}Problemas do uso}
    \begin{enumerate}
      \setcounter{enumi}{4}
      \item \textbf{Perda de habilidade}: atrofia pelo uso precoce.
      \item \textbf{Dados e privacidade}: o que você cola, você entrega.
      \item \textbf{Paradoxo da transparência}: declarar o uso pode ser penalizado.
    \end{enumerate}
  \end{columns}
  \vfill
  \small\textit{Os quatro primeiros nascem de como o modelo funciona e os três últimos nascem de como \textbf{nós} o usamos.}
\end{frame}
\begin{frame}{Viés: por que isso nos afeta em particular?}
  \begin{block}{O modelo aprendeu com o que estava disponível}
    E o que estava disponível é majoritariamente \textbf{em inglês}, do Norte global, de autores mais traduzidos e mais citados.
  \end{block}
  \pause
  \vfill
  \begin{itemize}
    \item A IA ``conhece'' muito melhor a literatura anglófona que a brasileira.
    \item O mesmo viés de dados de treinamento que documenta discriminação algorítmica contra grupos brasileiros marginalizados afeta também como a IA trata temas e fontes nacionais (Silva, 2022).
    \pause
    \item \alert{Para quem pesquisa o Brasil, isso não é detalhe. E sim, limite de fundo.}
  \end{itemize}
\end{frame}
\begin{frame}{Opacidade: por que não dá para abrir a caixa}
  \small
  \begin{block}{Três tipos de opacidade (Burrell, 2016)}
    \begin{itemize}
      \item \textbf{Sigilo intencional}: a empresa não revela como o modelo foi treinado.
      \item \textbf{Analfabetismo técnico}: poucos sabem ler código ou matemática o suficiente para entender, mesmo se revelado.
      \item \textbf{Opacidade de escala}: nem quem programou consegue explicar uma resposta específica, dado o tamanho do modelo.
    \end{itemize}
  \end{block}
  \pause
  \vfill
  \begin{alertblock}{}
    \centering Isso vale até para a Anthropic, a OpenAI e o Google. \textbf{Ninguém} consegue apontar, com certeza, por que o modelo escolheu esta palavra e não outra.
  \end{alertblock}
\end{frame}
\begin{frame}{Perda de habilidade e dados sensíveis}
  \textbf{Perda de habilidade}
  \begin{itemize}
    \item O paralelo é a calculadora e o GPS: ganhamos velocidade, perdemos o cálculo mental e o senso de direção.
    \item Numa formação de pesquisador, temos: a ferramenta está \textbf{ampliando} uma competência ou \textbf{impedindo} que ela se forme?
  \end{itemize}
  \pause
  \vspace{0.3cm}
  \textbf{Dados e privacidade}
  \begin{itemize}
    \item Colar um texto numa IA é \textbf{enviá-lo a uma empresa} (alimenta-la com nossos dados);
    \item Entrevistas, dados de campo e material sob sigilo, \alert{não vão para a IA} sem anonimização.\\ \textbf{Atenção à LGPD e ao comitê de ética.}
  \end{itemize}
\end{frame}
\begin{frame}{Irreprodutibilidade: o problema que atinge o método}
  \begin{alertblock}{A mecânica}
    O mesmo prompt, no mesmo modelo, produz respostas diferentes a cada execução. A variação não é erro, é parte de como o modelo gera texto.
  \end{alertblock}
  \pause
  \vfill
  \begin{itemize}
    \item Isso compromete diretamente a lógica de replicabilidade que é central ao método científico;
    \item Diferente de um erro humano replicável, aqui a variação é \textbf{estrutural} (Ouyang et al., 2025, \textit{ACM Transactions on Software Engineering and Methodology});
    \pause
    \item \alert{Implicação prática}, nunca cite ``a IA respondeu X'' como evidência fixa. \textbf{Registre data, versão do modelo e o prompt exato}. É o que sustenta a replicabilidade procedimental.
  \end{itemize}
\end{frame}
\begin{frame}{O paradoxo da transparência}
  \small
  \begin{alertblock}{O que pode acontecer}
    Declarar o uso de IA reduz a nota atribuída ao texto, mesmo quando o conteúdo é idêntico. Um experimento controlado mediu essa penalidade em 0,32 ponto, numa escala de 5, tanto para avaliadores humanos quanto para IA (Cheong et al., 2025).
  \end{alertblock}
  \pause
  \vfill
  \begin{itemize}
    \item A simples menção à IA já aciona rejeição, mesmo quando o uso foi legítimo (BaHammam, 2025);
    \item Isso cria um incentivo perverso: quem segue a regra ``declare todo uso'' pode ser penalizado por segui-la;
    \pause
    \item \alert{A resposta não é esconder!} É (1) reconhecer o problema e (2) cobrar mudança na cultura avaliativa, sem abandonar a transparência.
  \end{itemize}
\end{frame}
\begin{frame}{Detectores de IA}
  \small
  \begin{alertblock}{Uma consequência prática da opacidade e da irreprodutibilidade}
    Turnitin, GPTZero, Originality.ai e Copyleaks prometem identificar texto gerado por IA. Na prática, têm limitações documentadas que os tornam \textbf{inadequados como prova isolada}.
  \end{alertblock}
  \pause
  \begin{itemize}
    \item \textbf{Enviesados contra não nativos em inglês}. Liang et al. (2023), \textit{Patterns}, testaram sete detectores em redações TOEFL: \alert{61,3\% de falsos positivos} em textos de não nativos e quase zero em textos de nativos.
    \item \textbf{Fáceis de enganar}. Acurácia de até 98\% só em texto não editado; uma revisão leve, como a que fazemos neste curso, já derruba a detecção.
    \item \textbf{Sem validação robusta para o português} acadêmico.
  \end{itemize}
  \pause
  \vspace{0.05cm}
  \centering \textbf{Atenção}: o próprio fabricante do Turnitin orienta não tratar o percentual como prova definitiva.
\end{frame}
% ---------------- M4 ----------------
\subsection{M4. Ética, integridade e regulação}
\transicao{Módulo 4}{Ética, integridade e regulação}
\begin{frame}{O documento de referência da nossa área}
  \begin{exampleblock}{Leitura recomendada e usada na ANPOCS}
    \textbf{SAMPAIO; SABBATINI; LIMONGI} (2024). \textit{Diretrizes para o uso ético e responsável da Inteligência Artificial Generativa: um guia prático para pesquisadores}. São Paulo: Intercom.
  \end{exampleblock}
  \pause
  \begin{block}{O que ele diz, em três frases}
    \begin{itemize}
      \item \textbf{A IA não é coautora}: Quem responde pelo texto é você;
      \item \textbf{Declare o uso}: Diga qual ferramenta e para quê;
      \item \textbf{Verifique tudo}: especialmente as \textbf{referências}.
    \end{itemize}
  \end{block}
  \pause
  \centering\small São oito princípios no total. Estes três resolvem a maioria das dúvidas do dia a dia.
\end{frame}
\begin{frame}{O que as instituições decidiram (2025--2026)}
  \footnotesize
  \begin{columns}[t]
    \column{0.5\textwidth}
    \textbf{\color{azulUNI}Agências de fomento}
    \begin{itemize}
      \item \textbf{CNPq}: Portaria 2.664/2026, \alert{primeira norma federal obrigatória}. Autoria e responsabilidade integral do pesquisador;
      \item \textbf{CAPES}: sem portaria própria; Nota Técnica nº 3/2025 trata de integridade acadêmica e IA em teses, alinhada ao CNPq;
      \item \textbf{FAPESP}: sem documento específico, mas com regras de transparência em relatórios.
    \end{itemize}
    \column{0.5\textwidth}
    \textbf{\color{azulUNI}Universidades}
    \begin{itemize}
      \item \textbf{USP}: Portaria GR 9004/2026;
      \item \textbf{UFC}: Portaria 39/2025, pioneira e \textit{sem caráter punitivo};
      \item \textbf{Unifesp}: Resolução 17/2025 (CPGPq), declaração obrigatória em teses;
      \item \textbf{43\% das federais} (30 de 69) já têm protocolo ou constroem um (Andifes, 2026).
    \end{itemize}
  \end{columns}
  \pause
  \vspace{0.15cm}
  \begin{block}{}
    \centering A pergunta deixou de ser \textit{se} declarar o uso de IA. É \textbf{como} declarar.
  \end{block}
\end{frame}
\begin{frame}{USP, o ETDIA (Portaria GR 9004/2026)}
  \begin{exampleblock}{O que a portaria de fato cria}
    Em 10 de março de 2026, o Reitor criou, junto ao Gabinete do Reitor, o \textbf{Escritório de Transformação Digital e Inteligência Artificial (ETDIA)}.
  \end{exampleblock}
  \pause
  \begin{block}{Finalidade (Art. 2º)}
    Apoiar a USP na definição de diretrizes, na coordenação institucional e na promoção de boas práticas relacionadas à Transformação Digital e ao \textbf{uso ético e responsável da IA}.
  \end{block}
  \pause
  \centering\small É um escritório de governança, dirigido por um Coordenador indicado pelo Reitor, articulado com STI, EGIDA e o Escritório de Proteção de Dados.
\end{frame}
\begin{frame}{USP: o que o ETDIA vai fazer (Art. 3º)?}
  \small
  \begin{itemize}
    \item Propor diretrizes e políticas institucionais para TD e uso de IA no ensino, na pesquisa, na extensão e na gestão;
    \item Propor uma \textbf{ementa-base institucional em IA}, adaptável por área e curso;
    \item Apoiar a formação de docentes capacitadores em IA e competências digitais;
    \item Realizar benchmarking com universidades de referência, nacionais e internacionais.
  \end{itemize}
  \pause
  \vfill
  \begin{alertblock}{Atenção}
    A portaria cria o escritório e sua missão. Ela \textbf{não} estabelece, por si só, uma obrigação de declaração de uso de IA na pós-graduação. Regras desse tipo, se vierem, dependem de atos complementares da Reitoria (Art. 5º) ou de normas de cada unidade.
  \end{alertblock}
\end{frame}
\begin{frame}{FAPESP: o que está confirmado e o que falta verificar?}
  \begin{block}{O que é verificável hoje}
    A FAPESP mantém \textbf{Normas para Uso de Recursos e Prestação de Contas} e políticas para acesso abertos e propriedade intelectual, que tratam de autoria, integridade e transparência na pesquisa financiada.
  \end{block}
  \pause
  \begin{alertblock}{Lacuna}
    Nesta busca, \textbf{não localizei} uma norma específica e publicada da FAPESP sobre uso de IA generativa, comparável à Portaria 2.664/2026 do CNPq ou à do CNPq adotada pela CAPES.
  \end{alertblock}
\end{frame}
\begin{frame}{E no mundo e nos periódicos?}
  \small
  \begin{itemize}
    \item \textbf{União Europeia, AI Act, Artigo 4} (vigor desde fev/2025): o letramento em IA virou \alert{obrigação legal} de quem usa e oferece essas ferramentas, universidades incluídas;
    \pause
    \item \textbf{NTU Singapura}: toda tese leva uma declaração de originalidade dizendo como a IA foi usada (transparência sem punição);
    \pause
    \item \textbf{AJPS} (\textit{American Journal of Political Science}, política de ago/2026): o autor \textbf{pode} usar IA, mas declara em nota de rodapé; \alert{o parecerista não pode subir o manuscrito à IA}, só organizar as próprias ideias.
  \end{itemize}
  \pause
  \vspace{0.15cm}
  \begin{block}{Convergência internacional}
    Letramento $\bullet$ transparência $\bullet$ autoria humana $\bullet$ distinção entre uso do autor e do avaliador $\bullet$ declaração de uso.
  \end{block}
\end{frame}
\begin{frame}{Então: o que é permitido?}
  \footnotesize
  \begin{columns}[t]
    \column{0.33\textwidth}
    \textbf{\color{verdeCodigo}Em geral, aceito:}
    \begin{itemize}
      \item Revisar gramática e clareza
      \item Traduzir trechos
      \item Resumir texto \textit{que você leu}
      \item Sugerir palavras-chave
      \item Formatar referências \textit{que você tem}
      \item Apoiar código
    \end{itemize}
    \column{0.33\textwidth}
    \textbf{\color{laranjaDestaque}Aceito \textit{com} cuidado:}
    \begin{itemize}
      \item Estruturar argumentos
      \item Fichar textos
      \item Interpretar resultados
      \item Gerar \textit{abstract}
    \end{itemize}
    \vspace{0.1cm}
    \scriptsize\textit{Exige revisão crítica e, em geral, declaração.}
    \column{0.33\textwidth}
    \textbf{\color{laranjaDestaque}Vedado ou grave}
    \begin{itemize}
      \item IA como coautora
      \item Referências não verificadas
      \item Gerar resultados/dados
      \item Dados sensíveis sem anonimizar
      \item IA no manuscrito sob avaliação (parecer)
      \item Não declarar o uso
    \end{itemize}
  \end{columns}
  \pause
  \vspace{0.05cm}
  \begin{alertblock}{}
    \centering As cores não são lei. São apenas \textbf{grau de cuidado}. A norma do seu programa e do periódico sempre manda mais.
  \end{alertblock}
\end{frame}
\begin{frame}{Agora, a prática! Quatro casos}
  \small
  \begin{block}{Formato em 40 minutos}
    \begin{enumerate}
      \item Quatro casos concretos, projetados um a um.
      \item Em cada caso, \textbf{2 min} pensando sozinho e \textbf{8 min} discutindo em grupo.
      \item Registramos no quadro os critérios que a turma for construindo.
    \end{enumerate}
  \end{block}
  \pause
  \vspace{0.2cm}
  \begin{itemize}
    \item \textbf{Não existe gabarito}, mas um \textbf{critério defensável};
    \item Evitem os dois extremos: ``a IA resolve tudo'' e ``IA é fraude'';
    \pause
    \item \alert{No fim, queremos sair com uma lista de acordos práticos, nossa, do CEM.}
  \end{itemize}
\end{frame}
\begin{frame}{A revisão de literatura (Caso 1)}
  \begin{alertblock}{}
    \centering\large
    Um pós-graduando usa IA para redigir \textbf{toda} a revisão de literatura da dissertação. Depois lê tudo, corrige e verifica cada referência.
  \end{alertblock}
  \vfill
  \begin{itemize}
    \item É uso legítimo? Onde exatamente está a linha?
    \item Muda alguma coisa se ele \textbf{declarar} o uso?
    \item E se, em vez de ``toda a revisão'', fosse ``o primeiro rascunho''?
     \item Nesse caso, qual (is) acordo (s) você(s) sugere(m)?
  \end{itemize}
\end{frame}
\begin{frame}{As entrevistas (Caso 2)}
  \begin{alertblock}{}
    \centering\large
    Uma pesquisadora tem 40 entrevistas transcritas e cola todas numa IA para pedir uma \textbf{análise temática}.
  \end{alertblock}
  \vfill
  \begin{itemize}
    \item Quais problemas vocês enxergam aqui? (Pensem em \textbf{dois} tipos diferentes.)
    \item O que o comitê de ética diria? E a LGPD?
    \item Existe uma forma segura de fazer algo parecido?
         \item Nesse caso, qual (is) acordo (s) você(s) sugere(m)?
  \end{itemize}
\end{frame}
\begin{frame}{O parecer (Caso 3)}
  \begin{alertblock}{}
    \centering\large
    Um professor recebe um artigo para \textbf{parecer} de um periódico. Cola o manuscrito na IA e pede ajuda para avaliar.
  \end{alertblock}
  \vfill
  \begin{itemize}
    \item O que há de errado aqui que não havia nos casos anteriores?
    \item Por que a AJPS, e tantos outros periódicos, \textbf{vedam} isso ao parecerista?
    \item E se ele usasse a IA só para revisar o \textit{português} do parecer que ele mesmo escreveu?
         \item Nesse caso, qual (is) acordo (s) você(s) sugere(m)?
  \end{itemize}
\end{frame}
\begin{frame}{A aula (Caso 4)}
  \begin{alertblock}{}
    \centering\large
    Uma docente pede à IA que gere o \textbf{plano de aula} e a \textbf{lista de leituras} de uma disciplina nova e leva para a sala sem revisar.
  \end{alertblock}
  \vfill
  \begin{itemize}
    \item Qual o risco mais provável de acontecer? (Lembrem do Hands-on 1.)
    \item E se ela revisasse tudo? Continua sendo problema?
    \item Ela precisa contar aos alunos que usou IA?
         \item Nesse caso, qual (is) acordo (s) você(s) sugere(m)?
  \end{itemize}
\end{frame}

\begin{frame}{Fechando o debate}
  \begin{block}{O que ficou no quadro}
    \centering Os critérios que \textbf{vocês} construíram nos quatro casos.
  \end{block}
  \pause
  \vfill
  \begin{itemize}
    \item Voltando à pergunta do Módulo 2, a IA mudou \textit{o quê} ou \textit{o como} do nosso trabalho?
    \pause
    \item Lembrem do paradoxo da transparência e da irreprodutibilidade (Módulo 3). Eles não têm solução fácil e exigem julgamento constante, não uma regra decorada.
    \pause
    \item \alert{Por isso o critério importa mais que a regra decorada.}
  \end{itemize}
\end{frame}

\begin{frame}{Referências do Dia 1 (1/2), artigos e guias}
  \tiny
  \begin{itemize}
    \item ALKAISSI, H.; MCFARLANE, S. I. Artificial Hallucinations in ChatGPT: Implications in Scientific Writing. \textit{Cureus}, v. 15, e35179, 2023.
    \item AYDENIZ, M. Engineering Education at the Intersection of Knowledge, Work, and Power: Generative AI and Epistemic Authority. \textit{Sociology Compass}, 2026.
    \item BAHAMMAM, A. S. The Transparency Paradox: Why Researchers Avoid Disclosing AI Assistance in Scientific Writing. \textit{Nature and Science of Sleep}, 2025. DOI: 10.2147/NSS.S568375.
    \item BORIATI, D. From Research Tool to Epistemic Actor: Artificial Intelligence as Co-Producer of Social Knowledge. \textit{Societies}, v. 16, n. 6, p. 192, 2026.
    \item BURRELL, J. How the Machine Thinks: Understanding Opacity in Machine Learning Algorithms. \textit{Big Data \& Society}, v. 3, 2016. DOI: 10.1177/2053951715622512.
    \item CHEONG, I. et al. Penalizing Transparency? How AI Disclosure and Author Demographics Shape Human and AI Judgments About Writing. Preprint, arXiv:2507.01418, 2025.
    \item LIANG, W. et al. GPT detectors are biased against non-native English writers. \textit{Patterns}, v. 4, n. 7, 2023.
    \item OUYANG, S. et al. An Empirical Study of the Non-Determinism of ChatGPT in Code Generation. \textit{ACM Transactions on Software Engineering and Methodology}, 2025. DOI: 10.1145/3697010.
    \item PUNZIANO, G. Adaptive Epistemology: Embracing Generative AI as a Paradigm Shift in Social Science. \textit{Societies}, v. 15, n. 7, p. 205, 2025.
    \item SAMPAIO, R. C.; SABBATINI, M.; LIMONGI, R. Diretrizes para o uso ético e responsável da Inteligência Artificial Generativa. São Paulo: Intercom, 2024.
    \item SILVA, T. Racismo algorítmico: inteligência artificial e discriminação nas redes digitais. São Paulo: Edições Sesc, 2022.
  \end{itemize}
\end{frame}

\begin{frame}{Referências do Dia 1 (2/2), normas e legislação}
  \tiny
  \begin{itemize}
    \item AMERICAN JOURNAL OF POLITICAL SCIENCE. AJPS AI Policy. Atualizada em 11 ago. 2026. Disponível em: \url{https://ajps.org/ajps-ai-policy/}.
    \item ANDIFES. Levantamento sobre protocolos de IA nas universidades federais. 2026. \alert{[fonte de divulgação, verificar nota técnica original antes de citar em texto formal]}
    \item BRASIL. CAPES. Nota Técnica nº 3/2025, integridade acadêmica e uso de IA em dissertações e teses.
    \item BRASIL. CNPq. Portaria nº 2.664, de 6 de março de 2026. Institui a Política de Integridade na Atividade Científica. Disponível em: \url{gov.br/cnpq}.
    \item FAPESP. Norma específica sobre IA generativa não localizada nesta busca; ver Normas para Uso de Recursos e Prestação de Contas.
    \item NTU SINGAPORE. NTU Position on the Use of Generative Artificial Intelligence in Research, 2023.
    \item UFBA. Guia para Uso Ético e Responsável da Inteligência Artificial Generativa na UFBA, 2025--2026.
    \item UFC. Portaria nº 39/PRPPG/UFC, 2025.
    \item UFF. Guia para o Uso de Ferramentas de IA Generativa, out. 2025.
    \item UNIÃO EUROPEIA. AI Act, Artigo 4º e Artigo 50.
    \item UNIFESP. Resolução nº 17/2025 (CPGPq), diretrizes para uso de IA generativa na pós-graduação e na pesquisa.
    \item USP. Portaria GR nº 9004, de 10 de março de 2026. Cria o ETDIA.
  \end{itemize}
\end{frame}

\begin{frame}[plain]
  \centering
  \vspace*{2cm}
  {\Huge\color{azulUNI}\textbf{Fim do Dia 1}}\\[1cm]
  {\large Hoje: o que é, epistemologia, o que erra, o que as normas dizem.}\\[0.6cm]
  {\normalsize\color{laranjaDestaque} Amanhã: prompts, ensino, pesquisa, escrita acadêmica, agentes e \textit{skills}, e a síntese final.}
\end{frame}

% =====================================================================
\section{Dia 2}
\aberturadia{DIA 2}{Praticar e aplicar}{8h00 -- 12h30}
% =====================================================================

% ---------------- M5 ----------------
\subsection{M5. Escrevendo bons prompts}
\transicao{Módulo 5}{Escrevendo bons prompts: modelos, dicas e \textit{hands-on}}

\begin{frame}{Antes e depois}
  \begin{alertblock}{Prompt ruim}
    \centering\textit{``Me fale sobre capacidade estatal.''}
  \end{alertblock}
  \pause
  \begin{exampleblock}{Prompt bom}
    \small
    \textit{``Você é um professor de Ciência Política. Explique o conceito de capacidade estatal para alunos de graduação que nunca ouviram o termo. Estruture em: (1) definição em uma frase; (2) as três dimensões mais usadas na literatura; (3) um exemplo brasileiro concreto. Use no máximo 300 palavras, em português acadêmico acessível. Não cite referências, prefiro buscá-las eu mesmo.''}
  \end{exampleblock}
  \pause
  \centering\small O segundo diz \textbf{quem}, \textbf{o quê}, \textbf{como}, \textbf{quanto} e \textbf{o que não fazer}.
\end{frame}

\begin{frame}{A anatomia de um bom prompt}
  \centering
  \begin{tikzpicture}[node distance=0.35cm and 0.6cm]
    \tikzstyle{comp}=[rectangle, rounded corners, draw=azulUNI, very thick,
                      fill=cinzaSaida, text width=2.2cm, align=center, minimum height=1.5cm]
    \tikzstyle{seta}=[-{Latex[length=2mm]}, very thick, azulUNI]
    \node[comp] (papel) {\textbf{1. Papel}\\\tiny quem a IA deve ser};
    \node[comp, right=of papel] (tarefa) {\textbf{2. Tarefa}\\\tiny o que fazer};
    \node[comp, right=of tarefa] (contexto) {\textbf{3. Contexto}\\\tiny o que ela precisa saber};
    \node[comp, right=of contexto] (exemplos) {\textbf{4. Exemplos}\\\tiny como deve ficar};
    \node[comp, right=of exemplos] (formato) {\textbf{5. Formato}\\\tiny estrutura e tamanho};
    \draw[seta] (papel) -- (tarefa);
    \draw[seta] (tarefa) -- (contexto);
    \draw[seta] (contexto) -- (exemplos);
    \draw[seta] (exemplos) -- (formato);
  \end{tikzpicture}
  \pause
  \vfill
  \begin{block}{}
    \centering Nem todo prompt precisa dos cinco. Mas quando a resposta vem ruim, \\ quase sempre \textbf{falta um deles}.
  \end{block}
\end{frame}

\begin{frame}{Diagnóstico rápido, por que a resposta veio ruim}
  \small
  Antes de reescrever o prompt inteiro, identifique qual dos cinco elementos falhou. É quase sempre só um.
  \begin{center}
  \scriptsize
  \renewcommand{\arraystretch}{1.3}
  \begin{tabular}{p{2.6cm} p{9.6cm}}
    \toprule
    \textbf{Elemento que falta} & \textbf{Como isso aparece na resposta ruim} \\
    \midrule
    (1) Papel & Tom, nível ou vocabulário fora do esperado, técnico demais ou informal demais. \\
    (2) Tarefa & A IA foge do que foi pedido, responde outra coisa ou mistura vários pedidos. \\
    (3) Contexto & Resposta genérica, que serviria para qualquer disciplina ou qualquer público. \\
    (4) Exemplos & Conteúdo até correto, mas no formato, tom ou estilo que você não queria. \\
    (5) Formato & Estrutura, extensão ou forma de entrega fora do especificado. \\
    \bottomrule
  \end{tabular}
  \end{center}
\end{frame}

\begin{frame}[fragile]{Esqueleto para montar o seu prompt}
\begin{lstlisting}
[Papel] Você é [quem a IA deve ser: um professor, um revisor de
periódico, um estatístico...].

[Tarefa] [O que fazer, com um verbo de ação claro: resuma,
compare, revise, classifique, traduza...].

[Contexto] [O que ela precisa saber para fazer isso bem: público,
nível, disciplina, tamanho do material, restrições].

[Exemplos] [Como deve ficar o resultado, ou o que evitar: um
exemplo curto, um tom de referência, um contraexemplo].

[Formato] [Estrutura, extensão e forma de entrega: tópicos,
tabela, número de palavras, com ou sem referências].
\end{lstlisting}
\centering\small Copie, preencha os colchetes e ajuste. É o modelo 00 do kit de sobrevivência.
\end{frame}

\begin{frame}{Duas instruções que evitam retrabalho}
  \begin{exampleblock}{Feche sempre o prompt com}
    \small
    (1) \textit{``Pergunte antes de prosseguir sempre que houver dúvida''}; (2) \textit{``Só produza o arquivo final depois que eu confirmar com um OK''}.
  \end{exampleblock}
  \pause
  \vfill
  \begin{itemize}
    \item Essas duas linhas evitam retrabalho, alucinação por ambiguidade e edição fora do combinado.
    \pause
    \item \alert{Guarde os prompts que funcionam} num arquivo simples, organizado por tarefa. Da próxima vez, é copiar, ajustar e usar, não reescrever do zero.
  \end{itemize}
\end{frame}

\begin{frame}{Construa e refine o seu prompt (Hands-on 2)}
  \begin{block}{Escolha uma tarefa real sua}
    Um plano de aula, uma pergunta de pesquisa, um parágrafo a revisar, um e-mail formal. Escreva um prompt para ela usando os cinco elementos.
  \end{block}
  \vfill
  \begin{enumerate}
    \item Rode na sua IA e leia a resposta com espírito crítico.
    \item Troque o prompt (não a resposta) com o colega ao lado: falta algum dos cinco elementos?
    \item Ajuste só o elemento que faltou e rode de novo.
  \end{enumerate}
\end{frame}

\begin{frame}{Boas práticas para levar}
  \small
  \begin{columns}[t]
    \column{0.5\textwidth}
    \textbf{Ao pedir}
    \begin{itemize}
      \item Separe a instrução do texto (aspas ou marcadores).
      \item Peça o raciocínio em passos quando for complexo.
      \item Peça alternativas, como ``me dê 3 versões''.
      \item Peça autocrítica: ``o que há de fraco nisso?''
    \end{itemize}
    \column{0.5\textwidth}
    \textbf{Ao usar}
    \begin{itemize}
      \item Conversa nova para tarefa nova.
      \item Anexe só o trecho relevante.
      \item \textbf{Guarde os prompts que funcionaram.}
      \item Registre ferramenta, versão e data.
    \end{itemize}
  \end{columns}
  \pause
  \vspace{0.25cm}
  \begin{alertblock}{Regra de ouro nº 2}
    \centering Nenhuma saída entra no seu trabalho sem a sua revisão crítica.
  \end{alertblock}
\end{frame}

% ---------------- M6 ----------------
\subsection{M6. IA no ensino}
\transicao{Módulo 6}{IA no ensino: aulas e apresentações}

\begin{frame}{O que a IA faz bem na docência}
  \begin{columns}[t]
    \column{0.5\textwidth}
    \textbf{\color{verdeCodigo}Ajuda de verdade}
    \begin{itemize}
      \item Rascunhar plano de aula
      \item Gerar exercícios e questões
      \item Criar exemplos e analogias
      \item Adaptar linguagem ao nível da turma
      \item Elaborar rubricas de avaliação
      \item Resumir texto denso \textit{que você já leu}
    \end{itemize}
    \column{0.5\textwidth}
    \textbf{\color{laranjaDestaque}Cuidado}
    \begin{itemize}
      \item Lista de leituras (\alert{referências inventadas})
      \item Corrigir trabalho sem ler
      \item Conteúdo factual sem conferir
      \item Levar para a sala sem revisar
    \end{itemize}
  \end{columns}
  \pause
  \vfill
  \begin{block}{}
    \centering A IA rascunha. \textbf{Você ensina.}
  \end{block}
\end{frame}

\begin{frame}[fragile]{Modelo de prompt para plano de aula}
\begin{lstlisting}
Você é um professor universitário de Ciências Sociais, com
experiência em didática do ensino superior.
Elabore o plano de uma aula de 100 minutos para uma turma de
[graduação/pós] sobre: [INSERIR TEMA].
O plano deve conter:
1. Uma pergunta motivadora inicial.
2. Três blocos conceituais, com tempo estimado.
3. Uma atividade de discussão em grupo.
4. Síntese final.
Para cada bloco, indique o objetivo de aprendizagem e o
conteúdo central.
Não sugira bibliografia: eu mesmo seleciono as leituras.
Formato: tabela com a estrutura, seguida da descrição.
\end{lstlisting}
\centering\small\alert{Note a instrução negativa. Ela elimina o maior risco.}
\end{frame}

\begin{frame}[fragile]{Modelo de prompt para questões de avaliação}
\begin{lstlisting}
Você é professor(a) de [disciplina/nível]. Elabore 5 questões
sobre [TEMA]: 2 dissertativas, 2 de múltipla escolha com 4
alternativas plausíveis, 1 de análise de caso curto.

Para cada questão, indique: o objetivo de aprendizagem segundo a
taxonomia de Bloom, o tempo estimado de resposta e um gabarito
comentado. Não invente dados estatísticos ou fatos nos
enunciados: use apenas conceitos e categorias da literatura.
\end{lstlisting}
\centering\small Peça também uma tabela de especificação (questão x objetivo x peso sugerido).
\end{frame}

\begin{frame}{Apresentações com IA}
  \footnotesize
  \begin{block}{O que já existe}
    \begin{itemize}
      \item \textbf{Gamma}: gera a apresentação inteira, estrutura e imagens, a partir de um único prompt.
      \item \textbf{Claude} (Artifacts ou Cowork): monta slides a partir de um conteúdo que \textit{você} já validou.
    \end{itemize}
  \end{block}
  \pause
  \vspace{0.15cm}
  \begin{alertblock}{O risco}
    \begin{itemize}
      \item Nenhuma delas entrega uma apresentação pronta para a sala sem revisão.
      \item O ganho é de \textbf{tempo de montagem}, não de conteúdo. Texto, dados e exemplos exigem a mesma checagem de sempre.
    \end{itemize}
  \end{alertblock}
  \pause
  \vspace{0.1cm}
  \centering Use a IA para o rascunho visual. A aula e os exemplos continuam seus.
\end{frame}

\begin{frame}[fragile]{Modelo de prompt para apresentação}
\begin{lstlisting}
Você é um designer instrucional. A partir do plano de aula abaixo
(ou anexo), organize o conteúdo em uma sequência de slides para
uma apresentação de [N] minutos.

Para cada slide, indique: título, de 3 a 5 pontos-chave (frases
curtas, não parágrafos) e uma sugestão de elemento visual (gráfico,
esquema, imagem descritiva), quando fizer sentido.

Não invente dado, exemplo ou citação que não esteja no meu plano
de aula. Se um bloco for denso demais para um slide, sugira como
dividir em mais de um.
\end{lstlisting}
\centering\small Cole essa estrutura no Gamma, no Claude ou monte manualmente. A revisão de conteúdo continua sendo sua.
\end{frame}

\begin{frame}{E o outro lado: os alunos usam}
  \begin{block}{O fato}
    Eles já usam. A pergunta não é como impedir, é o que fazer.
  \end{block}
  \pause
  \vfill
  \begin{itemize}
    \item \textbf{Explicitar a regra} no plano de ensino, o que é permitido nesta disciplina e o que não é.
    \pause
    \item \textbf{Pedir declaração} de uso nos trabalhos, como a NTU faz com teses.
    \pause
    \item \textbf{Repensar a avaliação}. O que a IA faz em 10 segundos talvez não devesse valer nota. Defesa oral, processo, discussão em aula.
    \pause
    \item \alert{O teste da defesa oral}. Se o aluno não consegue defender o que entregou, o problema apareceu.
  \end{itemize}
\end{frame}

\begin{frame}{Checklist, antes de usar IA no ensino e na avaliação}
  \small
  \begin{itemize}
    \item[$\square$] Verifiquei conceitos, dados e exemplos gerados por IA antes de levar para a aula, contra alucinação?
    \item[$\square$] Deixei claro para a turma o que é aceito, aceito com cuidado ou vedado no uso de IA na disciplina?
    \item[$\square$] Testei se as questões geradas por IA têm resposta única, correta e defensável?
    \item[$\square$] Evitei usar IA para gerar parecer ou \textit{feedback} substantivo sobre um trabalho sem revisão integral minha?
    \item[$\square$] Atualizei minha própria declaração de uso de IA no plano de ensino ou no material do curso?
  \end{itemize}
\end{frame}

% ---------------- M7 ----------------
\subsection{M7. IA na pesquisa}
\transicao{Módulo 7}{IA na pesquisa: ferramentas de busca e experiências de IAs acadêmicas}

\begin{frame}{As três linhagens da IA na pesquisa}
  \small
  \begin{columns}[t]
    \column{0.33\textwidth}
    \begin{block}{\scriptsize Linhagem 1}
      \textbf{Recuperação com base real}\\[2pt]
      \scriptsize Semantic Scholar, Google Scholar, Connected Papers, ResearchRabbit, OpenAlex. Indexam e recomendam; não inventam.
    \end{block}
    \column{0.33\textwidth}
    \begin{block}{\scriptsize Linhagem 2}
      \textbf{Síntese generativa}\\[2pt]
      \scriptsize Claude, ChatGPT, Gemini. Trabalham sobre o material que \textit{você} fornece. Aqui mora o risco de alucinação.
    \end{block}
    \column{0.33\textwidth}
    \begin{block}{\scriptsize Linhagem 3}
      \textbf{Agentes de pesquisa autônoma}\\[2pt]
      \scriptsize \textit{Deep research}, Gemini Notebook sobre múltiplos documentos, Claude Code. Encadeiam passos com autonomia.
    \end{block}
  \end{columns}
  \pause
  \vfill
  \begin{alertblock}{A distinção que resolve tudo}
    Use a Linhagem 1 para \textbf{encontrar} literatura. Use a Linhagem 2 para \textbf{trabalhar} a literatura que você já tem. A Linhagem 3, agentes autônomos, é o tema inteiro do Módulo 9.
  \end{alertblock}
\end{frame}

\begin{frame}{Qual IA usar para qual tarefa}
  \scriptsize
  \renewcommand{\arraystretch}{1.15}
  \begin{tabular}{p{3.0cm} p{2.9cm} p{7.3cm}}
    \toprule
    \textbf{Tarefa} & \textbf{Categoria} & \textbf{Por quê} \\
    \midrule
    Ler e sintetizar um PDF longo & Claude (com arquivos) ou Gemini Notebook & Ancoradas no material carregado, reduzem a resposta fora da fonte. \\
    Buscar literatura com citação rastreável & Perplexity, Consensus, Elicit, OpenAlex & Retornam trechos e links da fonte junto da resposta. \\
    Mapear redes de citação & ResearchRabbit, Litmaps & Mostram a evolução da literatura e citações de um trabalho central. \\
    Organizar dezenas de textos & Gemini Notebook & Centraliza várias fontes num só espaço, com rastreabilidade. \\
    Apoiar código estatístico & Claude, ChatGPT (execução de código) & Geram e testam o código; a lógica estatística continua sendo sua. \\
    \bottomrule
  \end{tabular}
  \pause
  \vspace{0.2cm}
  \begin{exampleblock}{}
    \centering Ferramenta ancorada em fonte (você anexa o PDF) alucina menos que ferramenta de busca livre.
  \end{exampleblock}
\end{frame}

\begin{frame}{Simples ou complexo? O esforço certo do modelo}
  \small
  \begin{itemize}
    \item A maioria dos provedores oferece níveis de esforço: modelos rápidos para tarefas simples, modelos de raciocínio estendido para tarefas com várias etapas.
    \item O nível certo não é só questão de custo.
  \end{itemize}
  \pause
  \vfill
  \begin{alertblock}{Mais ``pensamento'' não é sempre melhor}
    \begin{itemize}
      \item Raciocínio estendido pode gastar centenas de tokens numa conta trivial sem ganho de acurácia, o \textit{overthinking} (Chen et al., 2025).
      \item Além de certo limiar de complexidade, o mesmo tipo de modelo tende a \textbf{piorar}, não melhorar (Shojaee et al., 2025, Apple Machine Learning Research).
    \end{itemize}
  \end{alertblock}
\end{frame}

\begin{frame}{Esforço baixo, médio e alto, na prática}
  \scriptsize
  \renewcommand{\arraystretch}{1.3}
  \begin{tabular}{p{4.1cm} p{4.1cm} p{4.1cm}}
    \toprule
    \textbf{Baixo esforço} & \textbf{Médio esforço} & \textbf{Alto esforço} \\
    \midrule
    Revisão ortográfica e gramatical & Fichamento padrão de um texto & Desenho da metodologia de uma pesquisa \\
    Traduzir um parágrafo curto & Formatar referências em ABNT & Sintetizar e comparar vários artigos \\
    Responder uma pergunta factual simples & Rascunho de e-mail formal & Depurar um script com erro não óbvio \\
    \bottomrule
  \end{tabular}
  \pause
  \vfill
  \begin{block}{Critério prático}
    \centering Resposta objetiva e fácil de checar? Baixo esforço resolve. Exige ponderar \textit{trade-offs} ou encadear várias etapas? Vale o modelo mais avançado.
  \end{block}
\end{frame}

\begin{frame}[fragile]{Modelo de prompt, organizar meu desenho de pesquisa}
\begin{lstlisting}
Você é um metodólogo de pesquisa em Ciências Sociais. Abaixo
estão minhas anotações sobre o desenho da pesquisa, mesmo
incompletas: pergunta, hipóteses, variáveis e ideias de método.
[COLE SUAS ANOTAÇÕES AQUI]

A partir SOMENTE do que eu escrevi, organize: (1) a pergunta de
pesquisa, sinalizando ambiguidade; (2) o tipo de desenho que já
indico (quanti, quali ou misto); (3) as variáveis já mencionadas,
organizadas; (4) lacunas que eu precise resolver antes de seguir.

Não proponha método, variável ou hipótese que eu não tenha
mencionado. Aponte a pergunta que falta responder, não a resposta.
\end{lstlisting}
\end{frame}

\begin{frame}[fragile]{Modelo de prompt, revisão sistemática (PRISMA)}
\begin{lstlisting}
Você é um pesquisador experiente em revisão sistemática,
seguindo o protocolo PRISMA 2020. Ajude-me a montar a
estratégia de busca sobre: [TEMA].

Produza: (1) 3 a 5 combinações de string de busca (operadores
booleanos) para bases como Scopus e Web of Science; (2)
critérios de inclusão e exclusão; (3) categorias a extrair de
cada artigo incluído; (4) um rascunho do fluxograma PRISMA.

Não simule ou invente resultado de busca: eu mesmo rodo as
strings nas bases e colo os resultados de volta.
\end{lstlisting}
\end{frame}

\begin{frame}{Três armadilhas na síntese de literatura}
  \begin{enumerate}
    \item \textbf{Achatamento}, o resumo apaga as tensões teóricas. Fica tudo parecendo consenso.
    \pause
    \item \textbf{Falsa familiaridade}, ler o resumo dá a \textit{sensação} de ter lido o texto. Não deu.
    \pause
    \item \textbf{Seleção enviesada}, ela devolve o que \textit{parece} relevante, não o que é central no campo.
  \end{enumerate}
  \pause
  \vfill
  \begin{block}{}
    \centering A IA é ótima para \textbf{organizar} o que você leu. \\ É péssima para \textbf{substituir} a leitura.
  \end{block}
\end{frame}

\begin{frame}{Dados sensíveis, a linha vermelha}
  \begin{alertblock}{Não vão para a IA generativa}
    Transcrições de entrevistas $\bullet$ dados de campo identificáveis $\bullet$ material sob sigilo $\bullet$ dados pessoais de terceiros $\bullet$ manuscrito alheio em avaliação.
  \end{alertblock}
  \pause
  \vfill
  \begin{itemize}
    \item Colar numa IA é \textbf{transferir a uma empresa}. Isso conflita com a LGPD e com o que você assinou no comitê de ética.
    \pause
    \item Alternativas possíveis, anonimização rigorosa \textit{antes}, modelos auto-hospedados, ou simplesmente \textbf{não usar}.
  \end{itemize}
\end{frame}

\begin{frame}{Fichamento comparativo (Hands-on 3)}
  \begin{exampleblock}{Cole ou anexe dois textos curtos sobre o mesmo tema}
    \small
    \textit{``Produza um fichamento comparativo dos dois textos abaixo:\\
    (1) uma tabela comparando: argumento central, método, principal achado, limitação;\\
    (2) uma síntese em 2 parágrafos: onde convergem, onde divergem.\\
    Não invente informação. Se algo não estiver explícito no texto, escreva `não informado'.''}
  \end{exampleblock}
  \vfill
  \begin{itemize}
    \item Depois: \textbf{compare com a sua leitura}. O que a IA captou bem? O que ela achatou?
    \item Perceba: a instrução ``não invente / marque não informado'' muda o resultado.
  \end{itemize}
\end{frame}

% ---------------- M8 ----------------
\subsection{M8. IA na escrita acadêmica}
\transicao{Módulo 8}{IA na escrita acadêmica: maximizar, produzir e reportar o uso}

\begin{frame}{Onde usar e onde não usar, na escrita}
  \begin{columns}[t]
    \column{0.48\textwidth}
    \textbf{\color{verdeCodigo}Defensável}
    \begin{itemize}
      \item Revisar clareza e gramática
      \item Apontar repetições
      \item Sugerir transições
      \item Reorganizar a estrutura
      \item Traduzir trechos
      \item Rascunhar \textit{abstract} do seu texto
    \end{itemize}
    \column{0.48\textwidth}
    \textbf{\color{laranjaDestaque}Problemático}
    \begin{itemize}
      \item Gerar o argumento original
      \item Inventar referências
      \item Substituir a leitura crítica
      \item Concluir sem dados
      \item Escrever a discussão por você
    \end{itemize}
  \end{columns}
  \pause
  \vfill
  \begin{alertblock}{O teste da defesa oral}
    \centering Se você não consegue defender oralmente cada parágrafo, o uso foi indevido.
  \end{alertblock}
\end{frame}

\begin{frame}{Seção por seção}
  \small
  \begin{tabular}{p{3.2cm} p{10.5cm}}
    \toprule
    \textbf{Seção} & \textbf{O que é razoável pedir} \\
    \midrule
    Introdução & Clareza e estrutura. \alert{A tese é sua.} \\
    Referencial teórico & Organizar e comparar o que \textit{você} selecionou. Verificar todas as referências. \\
    Metodologia & Só clareza da redação. \alert{Vários documentos restringem geração de metodologia.} \\
    Resultados & \textbf{Nada.} A IA não gera resultados. Pode ajudar a redigir o que você achou. \\
    Discussão & Fluência. \alert{A interpretação é o cerne autoral. Não terceirize.} \\
    Revisão final & Gramática, clareza, coesão. O uso mais seguro e aceito. \\
    \bottomrule
  \end{tabular}
\end{frame}

\begin{frame}[fragile]{Modelo de prompt, revisão de parágrafo mantendo a voz autoral}
\begin{lstlisting}
Sugira melhorias de clareza, coesão e precisão neste parágrafo,
MANTENDO minha voz e meu argumento. Não reescreva do zero.

Para cada mudança, apresente: trecho original, trecho sugerido,
tipo de mudança (clareza, coesão, precisão ou concisão) e uma
justificativa breve. Classifique cada uma como obrigatória (erro
real) ou opcional (preferência de estilo).

Preserve o registro acadêmico e a terminologia da área. Ao final,
ofereça 3 opções de reescrita do parágrafo completo (mais concisa,
mais formal, mais direta), para eu escolher ou combinar.
\end{lstlisting}
\end{frame}

\begin{frame}{Melhore um parágrafo seu (Hands-on 4)}
  \begin{exampleblock}{Cole um parágrafo real de um texto seu e use o modelo anterior}
    \small Repare na diferença entre \textbf{melhorar} e \textbf{substituir}.
  \end{exampleblock}
  \vfill
  \begin{itemize}
    \item Pedir a explicação de cada mudança mantém você no comando e ensina a escrever melhor.
    \item \alert{A IA trocou seu vocabulário técnico por termos genéricos? Isso é descaracterização.}
  \end{itemize}
\end{frame}

\begin{frame}[fragile]{Modelo de prompt, tradução de \textit{abstract} (PT para EN)}
\begin{lstlisting}
Traduza o abstract abaixo para inglês acadêmico (US English),
para submissão em periódico de [área]. Preserve a terminologia
técnica e o hedging proporcional à evidência (não fortaleça nem
enfraqueça as afirmações).

Monte um glossário PT-EN dos termos técnicos, para eu manter
consistência no manuscrito. Sinalize entre colchetes qualquer
trecho ambíguo no original. Informe a contagem de palavras,
considerando o limite de [N] do periódico.
\end{lstlisting}
\end{frame}

\begin{frame}[fragile]{Modelo de prompt, resposta estruturada a parecerista}
\begin{lstlisting}
Organize minha resposta ao comentário abaixo em: (1) comentário,
transcrito; (2) minha resposta, objetiva e respeitosa; (3)
alteração realizada no manuscrito; (4) localização (página/seção).

Classifique o comentário por tipo (metodológico, teórico, de
redação) e por prioridade. Não amenize nem exagere minha
concordância: mantenha o tom que eu indicar. Não invente
justificativa metodológica que eu não tenha fornecido.
\end{lstlisting}
\centering\small O kit de sobrevivência traz mais três modelos deste eixo: projeto para agência de fomento, carta de motivação (BEPE/PDSE) e parecer \textit{ad hoc} como revisor.
\end{frame}

\begin{frame}{Sua declaração de uso, modelo curto}
  \begin{exampleblock}{Adaptado de Sampaio, Sabbatini e Limongi (2024)}
    \itshape ``Este trabalho utilizou a ferramenta [nome, versão] de IA generativa para [tarefa específica]. Todas as saídas foram revisadas, validadas e editadas pelo autor, que se responsabiliza integralmente pelo conteúdo final.''
  \end{exampleblock}
  \pause
  \vfill
  \begin{itemize}
    \item \textbf{Onde}, nota de rodapé na primeira página, ou seção própria antes das referências.
    \item \textbf{O que dizer}, qual ferramenta, para quê, e que você responde pelo texto.
  \end{itemize}
\end{frame}

\begin{frame}{Sua declaração de uso, modelo detalhado (padrão AJPS)}
  \begin{exampleblock}{Segundo a AJPS AI Policy (2026)}
    \small\itshape ``Este trabalho utilizou [ferramenta, versão] em [datas]. Uso que moldou diretamente o conteúdo: [tarefa específica], com [grau de supervisão]. Uso que apenas apoiou a redação: [tarefa geral]. O autor revisou e verificou todo o conteúdo gerado por IA, incluindo citações e dados, e assume integral responsabilidade pelo manuscrito.''
  \end{exampleblock}
  \pause
  \vfill
  \centering\small Distinga usos que moldaram diretamente a pesquisa dos que só apoiaram a redação. Correção ortográfica pontual não precisa de declaração; qualquer uso que toque dado ou análise, sim. Na dúvida, declare.
\end{frame}

\begin{frame}{Escreva sua declaração de uso (Hands-on 5)}
  \begin{block}{}
    \centering Escolha um trabalho real que você tem em andamento e escreva agora a sua declaração, usando um dos dois modelos.
  \end{block}
  \vfill
  \begin{itemize}
    \item Guarde-a junto ao trabalho. É a primeira peça do seu próprio kit de sobrevivência.
    \item Voltamos a este ponto na síntese final do Módulo 10.
  \end{itemize}
\end{frame}

% ---------------- M9 ----------------
\subsection{M9. Desdobramentos contemporâneos}
\transicao{Módulo 9}{Desdobramentos contemporâneos: agentes, \textit{skills} e assistentes de código}

\begin{frame}{Além do chat, agentes de tarefa}
  \small
  \begin{block}{O que muda}
    \begin{itemize}
      \item Não sugere a próxima linha: recebe um objetivo, lê o material inteiro e planeja os passos.
      \item Edita, testa e corrige sozinho, em ciclos sucessivos, até cumprir o objetivo.
      \item Mais conhecidos: Claude Code (Anthropic) e Codex (OpenAI). Mesmo padrão no Gemini CLI e no Antigravity (Google).
    \end{itemize}
  \end{block}
  \pause
  \vfill
  \begin{alertblock}{Por que voltamos aqui}
    \centering Isso é a Linhagem 3 do Módulo 7, agentes de pesquisa autônoma, agora aplicada a código e a fluxos de trabalho inteiros.
  \end{alertblock}
\end{frame}

\begin{frame}{Um caso real, documentado}
  \begin{exampleblock}{Rakuten + Claude Code}
    A equipe de engenharia da Rakuten usou o Claude Code para implementar sozinho, em sete horas de trabalho autônomo, um método complexo numa biblioteca de 12,5 milhões de linhas de código, com \textbf{99,9\% de acurácia} (Anthropic, relato de caso).
  \end{exampleblock}
  \pause
  \vfill
  \begin{itemize}
    \item O ganho não é só velocidade: é encadear \textbf{muitas} etapas sem intervenção humana a cada uma.
    \item \alert{E a supervisão?} Concentrou-se na definição do objetivo e na checagem do resultado final, não em cada passo.
  \end{itemize}
\end{frame}

\begin{frame}{Cowork, agentes para quem não programa}
  \small
  \begin{block}{O que é}
    \begin{itemize}
      \item Versão da Anthropic pensada para quem não programa.
      \item Conecta pastas do seu computador e executa tarefas completas de arquivo e organização.
      \item Exemplos: gerar planilhas, documentos e apresentações; organizar os arquivos de um projeto de pesquisa. Com sua supervisão a cada etapa.
    \end{itemize}
  \end{block}
  \pause
  \vfill
  \begin{exampleblock}{}
    \begin{itemize}
      \item Use agentes para tarefas repetitivas de várias etapas.
      \item Para uma pergunta única e pontual, o chat comum ainda é mais rápido.
      \item A regra de ouro continua valendo: revise o resultado, nunca terceirize o juízo final.
    \end{itemize}
  \end{exampleblock}
\end{frame}

\begin{frame}{Duas filosofias de agente de código}
  \small
  \begin{columns}[t]
    \column{0.5\textwidth}
    \textbf{\color{azulUNI}Agent-first (Claude Code, Codex)}
    \begin{itemize}
      \item Você descreve o objetivo, o agente conduz.
      \item Roda no terminal, sem interface gráfica.
      \item Melhor para tarefas longas, com supervisão ao final.
    \end{itemize}
    \column{0.5\textwidth}
    \textbf{\color{azulUNI}IDE-first (Cursor)}
    \begin{itemize}
      \item Você conduz, a IA sugere e completa.
      \item Editor completo, revisão linha a linha.
      \item Melhor para quem quer aprovar cada mudança.
    \end{itemize}
  \end{columns}
  \pause
  \vfill
  \begin{block}{}
    \centering As duas categorias convergem rápido: em 2026, praticamente todas já oferecem modo agente, CLI e \textit{handoff} em nuvem. O critério de escolha é \textbf{quanto controle passo a passo você quer manter}.
  \end{block}
\end{frame}

\begin{frame}{\textit{Skills}, ensine a IA a repetir o seu fluxo}
  \small
  \begin{block}{O que é}
    \begin{itemize}
      \item Uma pasta com instruções, roteiros e, às vezes, pequenos scripts que a IA consulta.
      \item Executa uma tarefa recorrente sempre da mesma forma: formatar referências em ABNT, seguir o checklist de uma revisão sistemática, aplicar o seu estilo de escrita.
    \end{itemize}
  \end{block}
  \pause
  \vfill
  \begin{itemize}
    \item A plataforma carrega só um resumo de cada \textit{Skill} instalada, e abre as instruções completas quando a tarefa realmente precisa.
    \item Aciona-se digitando uma barra e o nome no início do prompt, por exemplo \texttt{/revisao-abnt}.
    \pause
    \item \alert{Regra prática}, se você está colando o mesmo prompt longo em conversas diferentes, esse é o sinal de transformá-lo numa \textit{Skill}.
  \end{itemize}
\end{frame}

\begin{frame}{A plataforma Claude, além do chat}
  \small
  \begin{description}
    \item[Projects] instruções e arquivos persistentes, para guardar suas preferências (``sempre em português acadêmico, ABNT, nunca invente referências'') sem repetir a cada conversa.
    \item[Artifacts] documentos, tabelas, slides e protótipos gerados e editáveis dentro da própria conversa, sem sair dela.
    \item[API] uso programático, para quem integra Claude a scripts e pipelines próprios (R, Python).
  \end{description}
  \pause
  \vfill
  \centering\small \textit{Skills} e agentes, vistos nos slides anteriores, também fazem parte dessa plataforma.
\end{frame}

\begin{frame}{Novidade, a marca d'água invisível}
  \small
  \begin{block}{O que é}
    Modelos Claude lançados a partir de 2 de agosto de 2026 embutem uma marca estatística imperceptível em todo texto gerado; imagens em SVG, PNG e JPG recebem metadados C2PA assinados. Resposta ao Artigo 50 do AI Act europeu.
  \end{block}
  \pause
  \vspace{0.1cm}
  \begin{alertblock}{O limite que importa}
    A marca \textbf{não prova autoria}. Um texto escrito por você e só revisado pela IA pode carregá-la. A \textbf{ausência} da marca também não prova o contrário, já que edição extensa e paráfrase reduzem a detecção.
  \end{alertblock}
  \pause
  \centering Detecção técnica não substitui declaração honesta. A regra ``declare todo uso'' continua sendo sua, não da ferramenta.
\end{frame}

\begin{frame}{Como pensar sobre isso}
  \begin{alertblock}{O alerta}
    ``Porque dá para automatizar'' não significa ``deve-se automatizar''. Em pesquisa, cada camada de automação é uma camada a mais onde o erro passa despercebido.
  \end{alertblock}
  \pause
  \vfill
  \begin{itemize}
    \item Domine primeiro o básico, \textbf{o prompt}. É o que dá 80\% do resultado.
    \item Automação faz sentido quando você já sabe exatamente o que quer e sabe conferir.
    \pause
    \item \textbf{Onde aprender}, documentação oficial das ferramentas, Anthropic Academy, tutoriais em vídeo. O kit final tem os links.
  \end{itemize}
\end{frame}

% ---------------- M10 ----------------
\subsection{M10. Síntese e encerramento}
\transicao{Módulo 10}{Síntese e encerramento}

\begin{frame}{O que vocês levam do Dia 2}
  \footnotesize
  \begin{enumerate}
    \item \textbf{A anatomia do prompt} e como diagnosticar por que uma resposta veio ruim (M5).
    \item \textbf{Modelos} para aula, avaliação e apresentação, e o teste da defesa oral (M6).
    \item \textbf{As linhagens da pesquisa}, qual IA usar em cada etapa e o esforço certo do modelo (M7).
    \item \textbf{Modelos} para revisão, tradução, resposta a parecerista, e sua declaração de uso (M8).
    \item \textbf{Agentes, \textit{skills} e assistentes de código}, o que já existe e por que exigem supervisão redobrada (M9).
    \item \textbf{Os checklists e as três regras}, para consultar sempre que a dúvida voltar.
  \end{enumerate}
\end{frame}

\begin{frame}{Sua declaração de uso, retomando}
  \begin{block}{}
    \centering Vimos no Módulo 8 dois modelos prontos, curto (Sampaio, Sabbatini e Limongi, 2024) e detalhado (padrão AJPS).
  \end{block}
  \pause
  \vfill
  \begin{itemize}
    \item Se você já escreveu a sua no Hands-on 5, guarde-a junto ao trabalho em andamento.
    \item Se ainda não escreveu, é o primeiro item do checklist a seguir.
  \end{itemize}
\end{frame}

\begin{frame}{Checklist geral, antes de usar IA em qualquer tarefa}
  \small
  \begin{itemize}
    \item[$\square$] Este texto tem entrevista, dado de campo ou material sob sigilo? Se sim, não cola sem anonimizar.
    \item[$\square$] Pedi para a IA não inventar referências e para dizer quando não sabe?
    \item[$\square$] Vou verificar cada referência e cada dado numa fonte primária antes de usar?
    \item[$\square$] Vou revisar tudo com o teste da defesa oral: consigo defender cada parágrafo?
    \item[$\square$] Escolhi o esforço do modelo, baixo, médio ou alto, de acordo com o tamanho da tarefa?
    \item[$\square$] Se essa tarefa se repete toda semana, já vale a pena transformá-la numa \textit{Skill}?
    \item[$\square$] Este uso precisa ser declarado, segundo a norma da minha instituição ou do periódico?
    \item[$\square$] Registrei ferramenta, versão, data e o prompt usado, para replicabilidade?
  \end{itemize}
\end{frame}

\begin{frame}{Checklists por perfil}
  \scriptsize
  \begin{columns}[t]
    \column{0.5\textwidth}
    \textbf{\color{verdeCodigo}Para estudantes, antes de entregar}
    \begin{itemize}
      \item[$\square$] Usei IA só nas etapas aceitas pela escala de cuidado (M4), não na geração de conteúdo substantivo?
      \item[$\square$] Anonimizei entrevistas e dados de campo antes de colar em qualquer chat?
      \item[$\square$] Conferi cada citação na fonte original?
      \item[$\square$] Incluí a declaração de uso no formato exigido?
      \item[$\square$] Consigo defender oralmente cada escolha e cada parágrafo?
    \end{itemize}
    \column{0.5\textwidth}
    \textbf{\color{laranjaDestaque}Para professores, antes de usar em sala}
    \begin{itemize}
      \item[$\square$] Verifiquei conceitos, dados e exemplos contra alucinação?
      \item[$\square$] Deixei claro para a turma o que é aceito, aceito com cuidado ou vedado?
      \item[$\square$] Testei se as questões geradas têm resposta única e defensável?
      \item[$\square$] Evitei gerar parecer substantivo sobre um trabalho sem revisão integral minha?
      \item[$\square$] Atualizei minha declaração de uso no plano de ensino?
    \end{itemize}
  \end{columns}
\end{frame}

\begin{frame}{Recursos para continuar}
  \small
  \begin{itemize}
    \item \textbf{Kit de Sobrevivência do Cientista Social com IA}, o guia de bolso deste curso: os 16 modelos de prompt completos, checklists por perfil e bibliografia comentada por trilha de estudo.
    \item \textbf{Repositório de normativas institucionais}, IEA/USP, \textit{Understanding AI}, com dezenas de guias de universidades brasileiras e estrangeiras, para achar a regra da sua própria instituição.\\
    \scriptsize\url{understandingai.iea.usp.br/guidelines-para-o-uso-da-ia-no-ensino}
    \item \small\textbf{Guia completo} de Sampaio, Sabbatini e Limongi (2024), Intercom.
    \item \textbf{AJPS AI Policy} como modelo de política de periódico na nossa área.
    \item \textbf{Anthropic Academy} e documentação oficial, para quem quiser ir além do Módulo 9.
  \end{itemize}
\end{frame}

\begin{frame}{Referências novas do Dia 2}
  \tiny
  \begin{itemize}
    \item CHEN, X. et al. Do not think that much for 2+3=? On the overthinking of o1-like LLMs. \textit{Transactions on Machine Learning Research}, 2025. Preprint: arXiv:2412.21187.
    \item SHOJAEE, P. et al. The illusion of thinking: understanding the strengths and limitations of reasoning models via the lens of problem complexity. Apple Machine Learning Research, 2025.
    \item ANTHROPIC. Skills explained: how Skills compares to prompts, Projects, MCP, and subagents. Disponível em: \url{https://claude.com/blog/skills-explained}. Acesso em: 13 ago. 2026.
    \item ANTHROPIC. Introduction to agentic coding; Claude Code overview e caso Rakuten. Disponíveis em: \url{https://claude.com/blog/introduction-to-agentic-coding} e \url{https://code.claude.com/docs/en/overview}. Acesso em: 13 ago. 2026. \alert{[relato institucional, não revisado por pares]}
    \item Comparativos de mercado sobre Claude Code, Cursor, Codex, Gemini CLI/Antigravity e sobre Gamma para apresentações (2026). \alert{[fontes jornalísticas e de mercado, sujeitas a mudança rápida; não peer-reviewed]}
  \end{itemize}
\end{frame}

\begin{frame}{As três regras}
  \vfill
  \begin{alertblock}{}
    \centering\Large
    Verificar toda referência\\[0.3cm]
    Declarar todo uso\\[0.3cm]
    Nunca terceirizar o juízo
  \end{alertblock}
  \pause
  \vfill
  \begin{block}{}
    \centering Se você esquecer todo o resto, estas três resolvem quase tudo.
  \end{block}
\end{frame}

\begin{frame}{Fechamento}
  \begin{block}{}
    \centering\large
    A IA é uma ferramenta útil e legítima, \textbf{quando usada com critério}.
  \end{block}
  \vfill
  \begin{itemize}
    \item Voltando à pergunta do Módulo 2, ela muda o \textit{como}. O \textit{quê} continua sendo nosso.
    \item A pergunta, o argumento, a interpretação e a responsabilidade não são delegáveis.
  \end{itemize}
  \vfill
  \centering\small\textit{Dúvidas, agora ou depois, estou à disposição.}
\end{frame}

\begin{frame}{Nota de transparência sobre este material}
  \begin{block}{Declaração de uso de IA}
    \small Este material foi produzido com apoio do \textbf{Claude (Anthropic), modelo Sonnet 5}, usado para estruturar o curso, redigir e revisar os slides, e verificar referências e normas citadas (agosto de 2026).
  \end{block}
  \pause
  \vfill
  \begin{itemize}
    \item Todo o conteúdo foi revisado, validado e editado pelo autor.
    \item O autor assume integral responsabilidade pelo conteúdo final.
  \end{itemize}
  \pause
  \vfill
  \centering\small\textit{A mesma regra que este curso ensina, aplicada aqui.}
\end{frame}

\begin{frame}[plain]
  \centering
  \vspace*{1cm}
  {\Huge\color{azulUNI}\textbf{Obrigado!}}\\[1.2cm]
  {\Large\textbf{Anderson Henrique}}\\[0.4cm]
  \small CEM/USP $\bullet$ INCT QualiGov\\
  \small Pesquisador Convidado IPEA\\
  \small Doutor em Ciência Política pela UFPE\\[1cm]
  \normalsize
  \begin{tabular}{rl}
    \textbf{E-mail} & andersonheri@gmail.com \\
                     & anderson.henrique@usp.br \\
    \textbf{Lattes} & \href{http://lattes.cnpq.br/7628365040600511}{lattes.cnpq.br/7628365040600511} \\
    \textbf{ORCID}  & 0000-0002-1842-2725 \\
  \end{tabular}
\end{frame}

\end{document}
